\section{Explaining Urban Social Structure}


The reviewed stylized facts about residential segregation, neighborhood change, and housing inequalities (in North America and Europe) provide many interesting explananda for urban research. Indeed, many theories have been proposed to understand these patterns and dynamics. In this section, I will briefly review existing theories. As many theories share features, I have grouped them into families of theories and discuss their merits and shortcomings with specific reference to a few theories of the respective families. Lastly, I summarize my critique of these theoretical models and present my reasons for developing a new one.

Explanations for neighborhood change usually describe triggers that distort a spatial equilibrium, whereafter the described patterns of neighborhood change follow until the city settles into a new spatial equilibrium \citep{galsterNatureNeighbourhood2001}. For example, gentrification is connected to the overall economic situation, state-funded infrastructure investment, changing tastes and preferences, population growth, and changes in income distributions \citep{leyAlternativeExplanationsInnerCity1986, zukinGentrificationCultureCapital1987}. I am not going to review potential triggers that can set off neighborhood change or reasons spatial equilibria become unstable. I am concerned with theories that describe the reasons and the processes that lead to residential segregation. Such theories should be able to produce the reviewed stylized facts about neighborhood change when confronted with an exogenous shock to the spatial equilibrium. Unfortunately, such connections have been rare: "One way in which scholarship on gentrification has been limited is that it has generally remained isolated from the study of persistent poverty and segregation" \citep[227]{hwangThingsChangeThings2016}.

%TODO mention institutions, zoning laws etc.


\subsection{Neighborhood Preferences}

There is a vast literature on theoretical models of residential segregation that explain segregation with households' preferences concerning their neighbors originating in urban ecology theory (for a review, see \cite{fossettEthnicPreferencesSocial2006}). Most of these models explain racial/ethnic segregation based on a preference for co-ethnic neighbors. As shown by \citet{schellingDynamicModelsSegregation1971}, such preferences lead to self-organization of segregation because individual moves create separate neighborhoods that realize the preferences of one group but not the other group. Moving into a neighborhood dominated by the outgroup does not improve residential satisfaction, even if the households prefer a more mixed neighborhood. Mixed neighborhoods are unstable, and high levels of segregation arise even with low levels of in-group favoritism.

However, some related models explain income segregation and more general patterns of social segregation as well. Such models assume that households prefer to live next to high-income neighbors (e.g. \cite{tieboutPureTheoryLocal1956, guerrieriEndogenousGentrificationHousing2013}) or that they value traits of their neighbors that are correlated with income (e.g. \cite{bayerEquilibriumModelSorting2004, fossettEthnicPreferencesSocial2006, benardWealthStatusBasedModel2007, yavasDissectingIncomeSegregation2019}). In these cases, different neighborhoods do not realize preferences for specific groups, but preferences are similar for all households, so the same neighborhoods are desirable for everyone. Segregation then arises because less affluent households can be outpriced. Households sort into the best neighborhoods they can afford. Such models can explain not only segregation by income but also the traits households have preferences over and those that correlate with income. Additionally, such models imply dynamics that closely mimic patterns of neighborhood change: when a neighborhood is becoming more attractive, poor residents are displaced by price increases and replaced by affluent in-movers. As more affluent neighbors have preferential characteristics, the neighborhood increases in attractiveness again. A self-reinforcing dynamic is created that continues until an increase in attractiveness no longer increases prices and the system reaches spatial equilibrium.

The main limitations of these models are that they fall short of the role of housing. Many models only consider preferences about neighbors (e.g., \cite{guerrieriEndogenousGentrificationHousing2013, yavasDissectingIncomeSegregation2019}). If the type and quality of the housing unit are unevenly distributed in space \citep{loufPatternsResidentialSegregation2016, oecdDividedCitiesUnderstanding2018, owensBuildingInequalityHousing2019, beaubrun-diantImpactPublicHousing2022} and preferences on housing characteristics are prioritized over preferences on the neighborhood \citep{chauCriticalReviewLiterature2003, bayerEquilibriumModelSorting2004, mummoloWhyPartisansNot2017, delucaNotJustLateral2020}, such models misrepresent the individual behavior leading to segregation. In the worst case, neighborhood preferences are not needed to generate segregation because of the uneven placement of more or less desirable housing units. Such models cannot explain the relationship between resident characteristics and housing quality. The models that consider preferences over housing units typically assume them to be exogenously given and fixed over time (e.g., \cite{bayerEquilibriumModelSorting2004, fossettEthnicPreferencesSocial2006, benardWealthStatusBasedModel2007}). While they can explain housing inequalities and better represent individual preferences, these models cannot explain the dynamics of housing quality. Changes in housing are a defining feature of neighborhood change \citep{zukinGentrificationCultureCapital1987, redfernWhatMakesGentrification2003} and therefore required in theories representing neighborhood change. Furthermore, the explanation of residential segregation by these models largely rests on the assumptions made about and does not explain the spatial distribution of housing quality.


\subsection{Amenities}

A typical family of explanations of both the spatial distribution of housing/neighborhood quality and socio-economic segregation is that more affluent households can outcompete poorer households for locations with the most amenities. Theories of this style have been influential in urban economics, particularly considering the relation of household income, housing unit size, and distance to the city center (see \cite{bruecknerLecturesUrbanEconomics2011}). In such models, housing is assumed to be cheaper with increasing distance to the central business district. Households can afford larger properties in suburbia at the expense of longer commutes. The choice of suburban living is made primarily by more wealthy families, causing an increase in household income and housing quality within a distance of the city center, while the poor live in small apartments in the city center. However, this concentric city model only performs well within the US \citep{bruecknerWhyCentralParis1999, quillianSocioeconomicSegregationLarge2016, gaigneWhoLivesWhere2022}. This model also does not represent neighborhood change well.

\citet{bruecknerWhyCentralParis1999} propose a theory based on more amenities than living space and commuting time. They consider how the (exogenously given) spatial distribution of amenities like rivers, hills, waterfronts, parks, and historical buildings affect whether affluent households reside in the city center or the suburbs. They explain the pattern that many US cities have poor city centers and affluent suburbs. In contrast, European cities tend to have rich centers and poor outskirts because of the historically grown concentration of amenities in the old European city centers that counteract the suburbanization forces described above. However, the theory can account for many income segregation patterns depending on the spatial distribution of amenities. \citet{leeNaturalAmenitiesNeighbourhood2018, gaigneWhoLivesWhere2022} are examples of more recent theories based on exogenous amenities. 

Amenity-based theories assume, often implicitly, that housing units of higher quality are built in places that are attractive to affluent households. However, only the classic concentric model directly addresses why different housing is built in different spaces: land prices decline with distance to the city center. So, these models can explain the joint patterns of people and infrastructure and show how they reinforce each other to create stable arrangements in space, even if the supply side is rarely given much agency. Nevertheless, exogenous (dis)amenities are neither necessary nor sufficient for residential segregation. For example, \citet{heblichEastSideStoryHistorical2021} show that historical pollution patterns in England created socio-economic segregation with the poor living in the most polluted areas. They also show that the geographical pattern remained stable even decades after the pollution stopped (see also \cite{ruttenauerEnvironmentalInequalityResidential2021}). There is an asymmetrical effect of amenities. While the placement of amenities can explain where specific groups segregate, it cannot explain why socio-economic segregation persists after amenities change. Some amenity-based theories include endogenous changes in amenities due to population changes, like restaurants and shops that cater to these demographics (e.g., \cite{bruecknerWhyCentralParis1999}) that would explain the inertia in segregation patterns after changes in amenities. However, they still need to pay more attention to changes in housing stock and supplier's agency.


\subsection{Reasons for a new Theory}

In summary, many theories of segregation and neighborhood change focus on demand without (explicitly) considering supply. As housing preferences are more important than neighborhood preferences \citep{chauCriticalReviewLiterature2003, bayerEquilibriumModelSorting2004, mummoloWhyPartisansNot2017, delucaNotJustLateral2020} and housing is unequally distributed in space \citep{loufPatternsResidentialSegregation2016, oecdDividedCitiesUnderstanding2018, owensBuildingInequalityHousing2019}, such theories misrepresent segregation dynamics and neighborhood formation. 

Another key objective of my proposed theoretical model is to bridge the gap between residential segregation models and neighborhood change, as well as residential mobility and housing inequality. Despite their inherent connection - they are all phenomena generated by the housing market - these research areas have remained mainly separate \citep{hwangThingsChangeThings2016}. Segregation represents stable patterns of neighborhood inequality, while neighborhood change involves rare transformations of single neighborhoods. The individual mobility behavior of residents creates segregation. The unequal placement of housing and households with different socio-economic characteristics in space implies that households with different characteristics live in housing units and neighborhoods of different quality.

By integrating these strands of literature, I aim to synthesize explanations of fundamental questions of urban research: Why do residents segregate into distinct neighborhoods? Why are neighborhoods so stable over time? Under what conditions do neighborhoods change? Why does neighborhood change follow typical patterns of invasion and succession of different socio-economic groups? Which residents move more frequently, and why? Who lives in housing units and neighborhoods of high and low quality? How much are households paying for housing?

Compared to existing theories, the model shall provide a more general explanation of multiple related phenomena. It shall be more realistic in its micro-foundations compared to models based on neighborhood preferences or amenities alone, acknowledging the agency of owners/landlords and the importance of the built environment in households' decision-making about where to move.